\documentclass[UTF8,fontset=fandol,zihao=-4]{ctexart}
\usepackage[a4paper,margin=25mm]{geometry}
\usepackage{amsmath,booktabs}
\usepackage[hidelinks]{hyperref}
\linespread{1.15}
\title{论文写作技能验证样稿}
\author{项目：面向局部社会信息的多智能体强化学习与合作演化研究}
\date{2026年9月15日}
\begin{document}
\maketitle
\noindent\textbf{阅读说明：}下节是依据已提供的工作1论文源文件撰写的正文样例，供判断行文风格与事实准确性；尚未并入学位论文。它不是已完成的整个方法章节，也不包含新增实验结果。
\section{邻居信息参与价值更新的动机与机制}

空间公共物品博弈中，个体的收益同时受到自身行为和周围个体行为的影响。一次合作之后获得较低收益，并不能单独说明合作在后续交互中仍然不利。邻居的行为及其获得的反馈，为个体判断当前策略提供了额外的参照。本研究据此在 Q-learning 的价值更新中引入邻居影响，考察相对表现信息如何改变个体的学习过程，以及这种改变能否促进群体合作。

模型将博弈交互范围与信息感知范围分别设置。每个智能体参与以自身及四个一阶邻居为中心的五个公共物品博弈组；感知参数 $M$ 则决定可用于状态构建和邻居评价的信息范围。当 $M=2$ 时，感知集合包含自身和十二个邻居，其中可作为参考对象的邻居有十二个。扩大感知范围改变了可获取的信息，并未扩大每个博弈组的成员范围。

个体维护自己的 Q 表，并依据感知范围内的平均声誉构造状态。动作执行后，模型将归一化的博弈收益与声誉变化加权组合，得到用于学习的综合奖励。随后先进行标准的时序差分更新，再计算邻居影响。因此，邻居比较所使用的是综合奖励，不能直接等同于未经处理的博弈收益。

邻居影响机制从候选邻居中选出综合奖励最高的个体。只有该邻居的奖励高于自身时，附加修正才生效：双方动作一致时，增加当前状态下已执行动作的 Q 值；动作不一致时，减小该 Q 值。修正幅度由正奖励差和控制参数 $\kappa$ 共同决定，并使用全局最大邻居奖励差进行归一化。完成更新后，智能体仍依据自己的 Q 表，通过 $\epsilon$-greedy 策略选择后续动作。邻居信息由此影响个体对动作价值的判断，而动作选择仍由个体自身的策略完成。

上述规则以邻居报告能够反映其真实表现为前提。若报告值被人为抬高，按最大综合奖励选择参考对象便可能反复选中异常发送者，进而影响价值更新。这一局限引出了后续研究：在保留个体 Q-learning 决策过程的基础上，进一步研究邻居消息的选择方式及其修正强度。该问题能否得到缓解，需要在明确的消息异常模型下，通过对照实验检验。

\clearpage
\section{工作2实验结论检查}
以下数值来自中期汇报，尚未用逐次运行结果核验。它们用于检查计算和论证范围，不构成实验复现。

\begin{center}
\begin{tabular}{lrr}
\toprule
方法 & 无故障合作率 & 10\%持续虚假节点\\
\midrule
单最佳邻居 & 0.959 & 0.138\\
E2E-LTNA-Q & 0.773 & 0.778\\
\bottomrule
\end{tabular}
\end{center}

\textbf{可由汇报值计算的量。}在异常条件下，新方法与基线的均值差约为 $0.778-0.138=0.640$，即64.0个百分点；无故障时则低18.6个百分点。基线从无故障到异常条件下降82.1个百分点，新方法的汇报均值上升0.5个百分点。最后一个差值不能据此解释为异常促进合作，也不能在没有逐次数据时称为统计显著。

\textbf{训练目标的解释。}若工作2沿用每人参与五组、每组合作成本为1的规则，且福利定义为所有个体真实原始收益之和，那么每组总收益为 $(r-1)n_C(g)$。每个合作者被五个组计入，因而
\[
 W(t)=\sum_i P_i(t)=5(r-1)N_C(t),\qquad
 \overline P(t)=5(r-1)f_C(t).
\]
固定 $r>1$、相同时间窗口和权重时，两者成正比。因此，“合作率未显式写入目标”不足以证明训练目标与合作率无关。应先核对工作2的实际目标函数；若包含其他代价、非线性变换、不同参数权重或不同评估窗口，需要另外分析。本次仅对上述公式做了小网格穷举校验，未运行工作2训练。

\textbf{统计与对照。}测试种子10至19提供十次运行，不能将每次运行中的900个节点或十万时间步作为独立重复。若比较采用配对设计，应保存共享的初始环境、异常节点及逐次结果；相同种子编号本身不足以说明环境随机量逐项匹配。还应区分一个固定控制器的测试波动与重新训练控制器产生的波动。

\textbf{优先补充的证据。}先核验现有结果，再在相同设置下比较小固定修正强度、仅学习选择、仅学习门控及完整模型，以判断提升分别来自哪里。各方法记录相同口径的长期福利、合作率、停止步数与是否达到停止条件。达到步数上限的运行也应保留；提前停止代表满足设定判据，不等于理论收敛。最终报告需给出逐次结果以及适合该实验设计的差值区间。

\section{验证范围与来源}
三个技能的入口配置已通过本机结构校验。写作样例核对来源为 \path{work1-paper/en_main.tex} 中的模型、Q-learning及邻居影响部分。统计部分检查了汇报均值运算，并对固定公共物品博弈规则下的福利关系进行精确算术校验。这些检查展示本次工作流程；没有进行有无技能的盲评，不能据此量化写作质量提升。

本次应用了 research-writing-skill、experimental-design 和 statistical-analysis。依照后两者的引用要求，在此记录工具来源：Kassis, T., Agarwal, V., He, Y., Patel, D., \& Brueckner, A. M. (2026). \textit{Scientific Agent Skills: A Library of Procedural Knowledge for Research Agents}. \href{https://doi.org/10.48550/arXiv.2609.00065}{arXiv:2609.00065}。记录已于2026年9月15日核对，当前修订日期为9月2日。
\end{document}
